% Style based on the supplied Waseem_Generic_AI_ML_GenAI_AgenticAI.tex.
% Python replaces template tokens. Edit layout here, content in JSON/portfolio JS.
\documentclass[9pt,letterpaper]{extarticle}
\usepackage[left=0.68in,right=0.68in,top=0.45in,bottom=0.45in]{geometry}
\usepackage[T1]{fontenc}
\input{glyphtounicode}
\pdfgentounicode=1
\usepackage[utf8]{inputenc}
\usepackage{lmodern}
\usepackage{microtype}
\usepackage{pifont}
\usepackage{xcolor}
\usepackage{enumitem}
\usepackage{titlesec}
\usepackage[hidelinks]{hyperref}
\usepackage{tabularx,array,parskip,ragged2e,needspace}
\pagestyle{empty}
\setlength{\parindent}{0pt}
\setlength{\parskip}{1pt}
\definecolor{accent}{HTML}{155B8A}
\titleformat{\section}{\large\bfseries\color{accent}}{}{0pt}{}[\vspace{0.5pt}\color{accent}\titlerule]
\titlespacing*{\section}{0pt}{3pt}{2pt}
% Compact subsections beneath the main accent-rule sections.
\titleformat{\subsection}{\normalsize\bfseries\color{accent}}{}{0pt}{}
\titlespacing*{\subsection}{0pt}{3pt}{1pt}
\setlist[itemize]{leftmargin=1.15em,itemsep=0.8pt,topsep=1pt,parsep=0pt,partopsep=0pt}
\newcommand{\role}[4]{%
  \Needspace{4\baselineskip}\textbf{#1}\hfill\textbf{#2}\\[-1pt]
  \textit{#3}\hfill\textit{#4}\par}
\newcommand{\entry}[2]{\Needspace{4\baselineskip}\textbf{#1}\hfill\textit{#2}\par}
\newcommand{\skillrow}[2]{\textbf{#1:} & #2\\}
\hypersetup{pdftitle={Waseem Raza, PhD - InterDigital Research Engineer REQ26-1144},pdfauthor={Waseem Raza, PhD}}
\begin{document}

\begin{center}
{\LARGE\bfseries Waseem Raza, PhD}{\normalsize\mdseries }\\[2pt]
{\large\bfseries Research Engineer \textbar{} AI-Enabled Wireless Systems \& Testbeds}\\[1pt]
{\small 5G RAN/Core \textbar{} 6G Research \textbar{} Digital Twins \textbar{} Agentic AI}\\[3pt]
{\footnotesize \ding{41}\enspace \href{mailto:waseem.placements@gmail.com}{waseem.placements@gmail.com} \textbar{} \href{mailto:waseemraza844@gmail.com}{waseemraza844@gmail.com}\\[1pt]New Jersey, United States \quad \ding{37}\enspace \href{tel:+17328039703}{+1 732-803-9703}\\[1pt]\href{https://www.linkedin.com/in/waseem-raza-phd-9a737858/}{LinkedIn} \textbar{} \href{https://github.com/WaseemRaza844}{GitHub} \textbar{} \href{https://scholar.google.com/citations?user=d58Drr0AAAAJ\&hl=en}{Google Scholar} \textbar{} \href{https://waseemraza844.github.io/portfolio/learning.html}{Coursera Courses} \textbar{} \href{https://waseemraza844.github.io/portfolio/}{Portfolio}}
\end{center}
\vspace{-5pt}
\Needspace{3\baselineskip}
\section{Professional Profile}
Research engineer and applied ML practitioner with a PhD in Electrical and Computer Engineering (2025) and 7+ years of combined research and engineering experience. Develop working AI-enabled wireless prototypes, Python-based simulations, experimental platforms and operational ML workflows; research spans 5G RAN intelligence, NR-NTN, CSI compression, propagation digital twins, outage management and learning-based optimization. Strong foundation in 5G RAN and core architecture, with interest in end-to-end 6G systems, edge/cloud and distributed AI. At AT\&T, improved automated root-cause classification performance from approximately 20\% to 80\%. Combine reproducible experiments, cross-layer reasoning, software development and peer-reviewed research.
\subsection{Technical Competencies}
\Needspace{2\baselineskip}\noindent{\small\textbf{Wireless and Cellular Systems:} 5G NR, LTE, NR-NTN, RAN, mobility, cell selection, CSI, radio propagation, network measurements, service assurance; 3GPP architecture knowledge: 5G Core, AMF, SMF, UPF, PDU sessions, registration, handover, QoS flows, policy control and edge computing concepts}\par
\Needspace{2\baselineskip}\noindent{\small\textbf{Testbeds, Simulation and Digital Twins:} TurboRAN experimental platform, SyntheticNET Python system-level simulator, 3GPP-aligned channel and mobility modeling, MATLAB, controlled measurements, scenario generation, performance metrics and repeatable experiments}\par
\Needspace{2\baselineskip}\noindent{\small\textbf{AI/ML:} PyTorch, TensorFlow/Keras, scikit-learn, LightGBM, classification, anomaly detection, GANs, autoencoders, transfer learning, domain adaptation, reinforcement learning, distributed-system optimization and robustness evaluation}\par
\Needspace{2\baselineskip}\noindent{\small\textbf{Agentic and Generative AI:} Python, LLM APIs, RAG, vector databases, embeddings, LangChain, LangGraph, LlamaIndex, tool use, multi-agent orchestration, CrewAI, AutoGen and evaluation; applied portfolio projects}\par
\Needspace{2\baselineskip}\noindent{\small\textbf{Research Software and Integration:} Python, C, C++, MATLAB, Linux, SQL, NumPy, pandas, Git, Docker, CI/CD, REST APIs, TCP/IP fundamentals, debugging, visualization, experiment automation and technical documentation}\par
\Needspace{2\baselineskip}\noindent{\small\textbf{Cloud and Distributed Workflows:} Azure Databricks, Spark/PySpark, Azure Data Factory, Synapse, Azure Functions, Azure DevOps, GitHub Actions; Kubernetes familiarity and edge/cloud architecture concepts}\par
\Needspace{2\baselineskip}\noindent{\small\textbf{Sensing and Experimental AI:} Inertial sensing, accelerometer/gyroscope time series, CNN/LSTM, feature extraction, context-aware classification, cross-user evaluation and uncertainty-aware inference}\par
\Needspace{4\baselineskip}\textbf{Research and Engineering Strengths}\par
\begin{itemize}
\item Translate research hypotheses into implementable algorithms, experimental software and measurable validation plans; document architecture, assumptions and limitations.
\item Investigate failure modes, distribution shifts and resource constraints through controlled experiments and systematic debugging.
\item Collaborate with wireless researchers, software engineers and domain stakeholders; mentor junior researchers and communicate findings through reports, presentations and publications.
\end{itemize}
\Needspace{3\baselineskip}
\section{Professional Experience}
\role{Machine Learning Engineer}{02/2025 - Present}{AT\&T, contractor via Innovatix Inc.}{Bedminster, NJ}
\begin{itemize}\interlinepenalty=10000
\item Own end-to-end development, testing, and operational integration of an AI-driven diagnostic system for critical E911 service-assurance workflows, improving automated root-cause classification performance from approximately 20\% to 80\%.
\item Implemented text-classification baselines using TF-IDF and LightGBM, and investigated domain-adaptive neural models and semi-supervised learning for heterogeneous operational domains.
\item Developed multimodal AI prototypes combining work-order text, location information, network context, and operational signals to support GenAI-assisted diagnostics for VoNR and 5G service assurance.
\item Developed model-validation, error-analysis and failure-mode assessment workflows; investigated autoencoder-based anomaly detection for PSAP outages and cross-domain robustness.
\item Partnered with network engineers, software teams, and public-safety stakeholders to translate operational requirements into validated AI workflows and production-oriented data assets.
\end{itemize}
\role{Data Scientist and AI/ML Engineer}{09/2024 - 12/2024}{Impac Exploration Services}{Weatherford, OK}
\begin{itemize}\interlinepenalty=10000
\item Built an LLM-enabled analytical pipeline for extracting and interpreting information from geochemical well data.
\item Developed Python and Azure Databricks, Azure Functions and Synapse workflows for ingestion, transformation, validation and repeatable analytical processing.
\item Automated ETL processing for large exploration datasets, reducing data-preparation time by 40\% and improving reproducibility.
\item Implemented cross-validation, data-quality, and consistency checks that improved reliability of analytical outputs by 20\%.
\item Collaborated with geologists to convert domain requirements and technical information into AI-ready datasets and actionable analytical outputs.
\end{itemize}
\role{Advanced-Degrees Intern}{06/2024 - 08/2024}{AT\&T Labs}{Austin, TX}
\begin{itemize}\interlinepenalty=10000
\item Analyzed 500K+ LTE/5G-NR diagnostic logs and applied ML methods to detect anomalies, reducing diagnostic time by 20\%.
\item Investigated 5G NR-NTN cell-selection behavior and improved decision algorithms by 15\% through data-driven analysis and validation.
\item Conducted structured analysis of large operational datasets and communicated technical findings to engineering teams supporting next-generation wireless systems.
\end{itemize}
\role{Advanced Wireless Technology Development Intern}{09/2022 - 01/2023}{Futurewei Technologies}{Schaumburg, IL}
\begin{itemize}\interlinepenalty=10000
\item Improved autoencoder-based channel-state-information compression by 25\% using transfer learning and systematic experimentation on large wireless datasets.
\item Conducted 50+ experiments to validate model performance, robustness, and generalization under changing data conditions.
\item Collaborated with research and engineering stakeholders to translate modeling results into recommendations for next-generation communication systems.
\end{itemize}
\role{Advanced-Degrees Intern}{06/2022 - 08/2022}{AT\&T Labs}{Bedminster, NJ}
\begin{itemize}\interlinepenalty=10000
\item Designed ML and statistical methods for analyzing users' timing in 5G standalone networks, improving analytical precision by 20\%.
\item Analyzed large-scale user trace data with SQL and visualization tools and communicated findings to senior technical leadership.
\item Converted raw operational traces into interpretable performance indicators to support diagnosis and network optimization.
\end{itemize}
\role{Graduate Research Assistant}{09/2019 - 07/2025}{AI4Networks, University of Oklahoma}{Tulsa, OK}
\begin{itemize}\interlinepenalty=10000
\item Designed reusable Python simulation and data-generation platforms for large-scale experimentation, including a 3GPP-compliant network simulator and an AI-enabled 5G testbed.
\item Built transfer-learning, adversarial-learning, GAN, autoencoder, and reinforcement-learning pipelines for sparse-data and cross-domain modeling.
\item Developed robustness evaluation frameworks for models operating under distribution shifts, missing data, and limited labeled samples; led 100+ experiments from problem formulation through validation.
\item Developed scalable self-healing and outage-management frameworks that improved modeled network resilience by 30\%.
\item Published peer-reviewed research, presented technical results internationally, and mentored junior researchers in ML implementation, experimental design, and technical writing.
\end{itemize}
\Needspace{3\baselineskip}
\section{Education}
\entry{PhD, Electrical and Computer Engineering}{09/2019 - 07/2025}
University of Oklahoma, Norman, Oklahoma, USA\par
\entry{MSc, Telecommunication Engineering}{08/2014 - 08/2016}
University of Engineering and Technology, Taxila, Pakistan\par
\entry{BSc, Telecommunication Engineering}{11/2010 - 06/2014}
University of Engineering and Technology, Taxila, Pakistan\par
\Needspace{3\baselineskip}
\section{Selected Projects \& Research}
\Needspace{2\baselineskip}
\subsection{Projects \& Applied Research}
\entry{SyntheticNET: 3GPP-Compliant Network Simulation}{Research Project \textbar{} 2020-2024}
Developed a Python-based, 3GPP-aligned simulation environment for realistic mobility, propagation, and cellular-network experimentation when measured data is sparse or costly. Supported controlled data generation and repeatable studies in propagation modeling, network optimization, digital twins, and resilient wireless AI.\par
{\small\textit{Tools \& Methods:} Python, 3GPP, Wireless Simulation, Mobility Modeling, Propagation Modeling, Synthetic Data, Experiment Design, Digital Twins, Network Analytics\par}
\vspace{1pt}
\entry{TurboRAN Experimental Platform}{Research Infrastructure \textbar{} 2020-2024}
Built and used a multi-band, multi-tier experimental platform for evaluating AI-enabled radio-access-network methods under controlled network conditions. Supported repeatable measurement collection, testbed experimentation, RAN optimization studies, and validation of emerging wireless-network algorithms.\par
{\small\textit{Tools \& Methods:} 5G Testbeds, RAN, Network Measurement, Experiment Design, Wireless Prototyping, Multi-Band Systems, Performance Evaluation, RAN Optimization\par}
\vspace{1pt}
\entry{Propagation Models \& Digital Twins}{Research Project \textbar{} 2024}
Developed robust data-driven radio-propagation models for wireless digital twins using domain-informed generative learning and scarce measured data. Combined 3GPP-informed pretraining, GAN-based augmentation, domain adaptation, and statistical validation to improve resilience under distribution shift.\par
{\small\textit{Tools \& Methods:} Propagation Modeling, Digital Twins, GANs, Domain Adaptation, 3GPP Models, Synthetic Data, Statistical Validation, Wireless AI, Distribution Shift\par}
\vspace{1pt}
\entry{\href{https://github.com/WaseemRaza844/agentic-ai-projects-by-was}{Agentic AI and RAG Applications}}{Portfolio Project \textbar{} 2026}
Built grounded GenAI workflows that combine document retrieval, vector search, stateful orchestration, tool-using agents, memory, evaluation, and multi-agent handoffs. Emphasized controllable execution, traceable context flow, and reliable responses for multi-step tasks rather than single unconstrained model calls.\par
{\small\textit{Tools \& Methods:} Python, LangGraph, LangChain, CrewAI, AutoGen, RAG, Vector Databases, Embeddings, Tool Calling, Agent Memory, Evaluation\par}
\vspace{1pt}
\entry{Wearable Intelligence for Health Monitoring}{Research Project \textbar{} 2024-2026}
Developed deep-learning workflows for recognizing medication-taking and dietary activities from smartwatch accelerometer and gyroscope signals in free-living settings. Evaluated windowing, feature extraction, CNN/LSTM-style sequence models, cross-user generalization, and recall-oriented performance for health-monitoring applications.\par
{\small\textit{Tools \& Methods:} Wearable AI, Smartwatch IMU, Time Series, Deep Learning, CNN/LSTM, Feature Engineering, Cross-User Evaluation, Activity Recognition, Healthcare AI\par}
\vspace{1pt}
\entry{5G + Edge + IoT + AI Architecture Analysis}{Certification Project \textbar{} 2026}
Evaluated how 5G connectivity, edge computing, IoT, and AI interact in service architectures and enterprise deployment scenarios. Analyzed use-case requirements, edge placement, IoT data flows, AI processing layers, and business/operational tradeoffs for connected intelligent systems.\par
{\small\textit{Tools \& Methods:} 5G, Edge Computing, IoT, Artificial Intelligence, Architecture Analysis, Use-Case Evaluation, Data Flows, Deployment Tradeoffs\par}
\vspace{1pt}
\entry{AI-Ready Lakehouse Pipeline}{Portfolio Project \textbar{} 2024}
Designed an AI-ready lakehouse pipeline that prepares structured and semi-structured data for analytics and downstream machine-learning workloads. Implemented repeatable ingestion, cleansing, transformation, validation, and scalable processing patterns suitable for reproducible cloud data workflows.\par
{\small\textit{Tools \& Methods:} Azure Databricks, Apache Spark, PySpark, ETL/ELT, Data Ingestion, Schema Transformation, Data Validation, Feature Preparation, Cloud Data Engineering\par}
\vspace{1pt}
\Needspace{2\baselineskip}
\subsection{Publications}
\entry{Real-World Medication Adherence Monitoring Using Hybrid Neural Networks and Smartwatches}{2026}
\textit{HealthINF \textbar{} Accepted}\par
\entry{\href{https://doi.org/10.3390/s25247651}{On Multi-Parameter Optimization and Proactive Reliability in 5G and Beyond Cellular Networks}}{2025}
\textit{Sensors 2025}\par
\entry{\href{https://doi.org/10.1109/PIMRC59610.2024.10817403}{An AI-Driven Framework for Enhancing Resilience in Propagation Models to Enable Digital Twin}}{2024}
\textit{IEEE PIMRC}\par
\entry{Data-Driven Intelligent Outage Management for High Shadowing Environments in 5G\&B Networks}{2024}
\textit{IEEE SmartNets}\par
\entry{\href{https://arxiv.org/abs/2203.15899}{Toward a Hybrid RF/Optical Lunar Communication System (LunarComm)}}{2022}
\textit{IEEE Network}\par
\entry{Towards Positioning Error Impact Characterization and Minimization in User-Centric RAN}{2022}
\textit{IEEE WCNC}\par
\entry{\href{https://doi.org/10.1016/j.aeue.2020.153244}{Unification of RF Energy Harvesting Schemes under Mixed Rayleigh-Rician Fading Channels}}{2020}
\textit{AEÜ - International Journal of Electronics and Communications}\par
\Needspace{3\baselineskip}
\section{Certifications \& Professional Development}
\Needspace{3\baselineskip}\noindent\textbf{1. 4G Network Fundamentals} {\small --- \textcolor{blue}{\underline{\href{https://www.coursera.org/account/accomplishments/verify/DLW1YN474VZJ}{Institut Mines-Télécom / Coursera}}} \textbar{} Completed \textbar{} August 2, 2026 \textbar{} \textcolor{blue}{\underline{\href{https://waseemraza844.github.io/portfolio/certificates/wireless/4g-network-fundamentals.pdf}{Certificate}}}}\par
\textbf{Focus:} Covers LTE/4G architecture from radio access through core-network functions, including cellular-network evolution, mobility, resource management, and end-to-end service operation. Emphasizes how major network components interact to deliver reliable broadband connectivity.\par
{\small\textbf{Tools \& Skills:} LTE/4G architecture, RAN and EPC concepts, mobility and handover, radio resource management, cellular networking, network-service fundamentals\par}
\vspace{2pt}
\Needspace{3\baselineskip}\noindent\textbf{2. Business Considerations for 5G with Edge, IoT, and AI} {\small --- \textcolor{blue}{\underline{\href{https://www.coursera.org/account/accomplishments/verify/SAUEVQXV5W71}{EDUCBA / Coursera}}} \textbar{} Completed \textbar{} February 18, 2026 \textbar{} \textcolor{blue}{\underline{\href{https://waseemraza844.github.io/portfolio/certificates/wireless/bus-cons-5g.pdf}{Certificate}}}}\par
\textbf{Focus:} Examines how 5G capabilities combine with edge computing, IoT, and AI to enable new enterprise and telecommunications services. Emphasizes use-case evaluation, business value, deployment considerations, and technology tradeoffs across connected intelligent systems.\par
{\small\textbf{Tools \& Skills:} 5G business cases, edge computing, IoT, AI integration, use-case analysis, deployment tradeoffs, telecom strategy\par}
\vspace{2pt}
\Needspace{3\baselineskip}\noindent\textbf{3. Agentic AI with LangChain and LangGraph} {\small --- \textcolor{blue}{\underline{\href{https://www.coursera.org/account/accomplishments/verify/L4EY4SLGCPUE}{Coursera}}} \textbar{} Completed \textbar{} February 9, 2026}\par
\textbf{Focus:} Covers design of stateful agentic applications using LangChain and LangGraph, including graph-based orchestration, tool use, memory and state, routing, and multi-step workflows. Emphasizes controllable agent execution and composition of reliable LLM-driven processes.\par
{\small\textbf{Tools \& Skills:} LangChain, LangGraph, tool-calling, state and memory, workflow orchestration, routing, agent design\par}
\vspace{2pt}
\Needspace{3\baselineskip}\noindent\textbf{4. AI Foundations \& Industry Overview for Telecommunication} {\small --- \textcolor{blue}{\underline{\href{https://www.coursera.org/account/accomplishments/verify/3900EVOH4I1Y}{Coursera}}} \textbar{} Completed \textbar{} June 19, 2026 \textbar{} \textcolor{blue}{\underline{\href{https://waseemraza844.github.io/portfolio/certificates/wireless/c1-ai-found-indus-overview.pdf}{Certificate}}}}\par
\textbf{Focus:} Introduces AI and machine-learning concepts in the telecommunications domain, connecting network data with automation, optimization, predictive analytics, and operational decision support. Reviews how AI can support network planning, service assurance, and customer-facing workflows.\par
{\small\textbf{Tools \& Skills:} telecom AI/ML use cases, network analytics, automation, optimization, predictive modeling concepts, operational intelligence\par}
\vspace{2pt}
\Needspace{3\baselineskip}\noindent\textbf{5. Introduction to Deep Learning \& Neural Networks with Keras} {\small --- \textcolor{blue}{\underline{\href{https://www.coursera.org/account/accomplishments/verify/Y6NP4GWLS4DZ}{IBM / Coursera}}} \textbar{} Completed \textbar{} August 12, 2026}\par
\textbf{Focus:} Introduces deep-learning and neural-network fundamentals using Keras, covering network architecture, training, validation, and model improvement. Connects concepts such as activation functions and optimization with practical neural-network implementation and evaluation.\par
{\small\textbf{Tools \& Skills:} Keras, deep learning, neural networks, model training, validation, activation functions, optimization, model evaluation\par}
\vspace{2pt}
\Needspace{3\baselineskip}\noindent\textbf{6. Machine Learning} {\small --- \textcolor{blue}{\underline{\href{https://www.coursera.org/account/accomplishments/verify/VU3Y5TBLVLZT}{Stanford Online / Coursera}}} \textbar{} Completed \textbar{} March 29, 2020}\par
\textbf{Focus:} Covers foundational machine-learning methods for supervised and unsupervised problems, including regression, classification, model evaluation, regularization, clustering, and practical model-selection considerations. Emphasizes translating data into predictive models while controlling generalization error.\par
{\small\textbf{Tools \& Skills:} supervised learning, regression, classification, regularization, clustering, model evaluation, bias-variance tradeoffs, machine-learning workflows\par}
\vspace{2pt}
\end{document}
